Betrachten Sie die Turing-Maschine über dem Bandalphabet
$\Gamma=\{\blank,\texttt{0},\texttt{1}\}$ mit dem Zustandsdiagramm
\begin{center}
\def\l{3}
\begin{tikzpicture}[>=latex,thick]
\coordinate (q0) at (0,0);
\coordinate (q1) at (\l,0);
\coordinate (q2) at ({2*\l},0);
\coordinate (q3) at ({3*\l},0);
\coordinate (q4) at ({1.5*\l},{0.5*\l});
\coordinate (qa) at ({0.5*\l},{-0.5*\l});
\coordinate (qr) at ({2.5*\l},{-0.5*\l});
\node at (q0) {$q_0$};
\node at (q1) {$q_1$};
\node at (q2) {$q_2$};
\node at (q3) {$q_3$};
\node at (q4) {$q_4$};
\node at (qa) {$q_{\text{accept}}$};
\node at (qr) {$q_{\text{reject}}$};
\draw (q0) circle[radius=0.35cm];
\draw (q1) circle[radius=0.35cm];
\draw (q2) circle[radius=0.35cm];
\draw (q3) circle[radius=0.35cm];
\draw (q4) circle[radius=0.35cm];
\draw (qa) ellipse(0.65cm and 0.35cm);
\draw (qa) ellipse(0.60cm and 0.30cm);
\draw (qr) ellipse(0.65cm and 0.35cm);
\draw (qr) ellipse(0.60cm and 0.30cm);
\draw[->,shorten >= 0.35cm] ($(q0)+(-1.5,0)$) -- (q0);

\draw[->,shorten >= 0.35cm,shorten <= 0.35cm]
	(q0) to[out=60,in=120,distance=1cm] (q0);
\node at ($(q0)+(0,1)$) {$\scriptstyle \texttt{0}\to\texttt{0},R$};

\draw[->,shorten >= 0.35cm,shorten <= 0.35cm] (q0) -- (q1);
\node at ($0.5*(q0)+0.5*(q1)$) [above]
	{$\scriptstyle\texttt{1}\to\texttt{1},R$};
\draw[->,shorten >= 0.35cm,shorten <= 0.35cm]
	(q1) to[out=60,in=120,distance=1cm] (q1);
\node at ($(q1)+(0,1)$) {$\scriptstyle \texttt{1}\to\texttt{1},R$};

\draw[->,shorten >= 0.35cm,shorten <= 0.35cm] (q1) -- (q2);
\node at ($0.5*(q1)+0.5*(q2)$) [above]
	{$\scriptstyle\texttt{0}\to\texttt{0},L$};

\draw[->,shorten >= 0.35cm,shorten <= 0.35cm] (q2) -- (q3);
\node at ($0.5*(q2)+0.5*(q3)$) [above]
	{$\scriptstyle\texttt{1}\to\texttt{0},R$};

\draw[->,shorten >= 0.35cm,shorten <= 0.35cm] (q3) to[out=135,in=0] (q4);
\node at ($(q2)+(0.5,1.5)$) [above] {$\scriptstyle\texttt{0}\to\texttt{1},L$};
\draw[->,shorten >= 0.35cm,shorten <= 0.35cm]
	(q4) to[out=60,in=120,distance=1cm] (q4);
\node at ($(q4)+(0,1.1)$)
	{$\texttt{0}\to\texttt{0},L\atop\texttt{1}\to\texttt{1},L$};

\draw[->,shorten >= 0.35cm,shorten <= 0.35cm]
	(q4) to[out=180,in=45] (q0);
\node at ($(q1)+(-0.5,1.5)$) [above] {$\scriptstyle\blank\to\blank,R$};

\draw[->,shorten >= 0.65cm,shorten <= 0.35cm]
	(q0) to[out=-90,in=180] (qa);
\draw[->,shorten >= 0.65cm,shorten <= 0.35cm]
	(q1) to[out=-90,in=0] (qa);
\node at ($(q0)+(0.2,{-0.35*\l})$) [left] {$\scriptstyle\blank\to\blank,R$};
\node at ($(q1)+(-0.2,{-0.35*\l})$) [right] {$\scriptstyle\blank\to\blank,R$};

\draw[->,shorten >= 0.65cm,shorten <= 0.35cm]
	(q2) to[out=-90,in=180] (qr);
\draw[->,shorten >= 0.65cm,shorten <= 0.35cm]
	(q3) to[out=-90,in=0] (qr);
\node at ($(q2)+(0.2,{-0.35*\l})$) [left] {$
\texttt{0}\to\texttt{0},R
\atop
\ifthenelse{\boolean{loesungen}}{
{\color{red}\blank\to\blank,R}
}{}
$};
\node at ($(q3)+(-0.2,{-0.35*\l})$) [right] {$
\texttt{1}\to\texttt{1},R
\atop
\ifthenelse{\boolean{loesungen}}{
{\color{red}\blank\to\blank,R}
}{}$};

\end{tikzpicture}
\end{center}
\begin{teilaufgaben}
\item
Im Zustandsdiagramm fehlen ein paar Übergänge für eine deterministische
Turing-Maschine.
Ergänzen Sie wo nötig Übergänge, die den Bandinhalt nicht verändern und
nach $q_{\text{reject}}$ führen.
\item
Bestimmen Sie die Berechnungsgeschichte für die Verarbeitung des Wortes
\texttt{1010}.
\item
Was steht nach der Verarbeitung des Wortes \texttt{0101} auf dem Band?
\item
Gibt es Wörter aus $\Sigma^*$ mit $\Sigma = \{\texttt{0},\texttt{1}\}$, die
von dieser Turing-Maschine nicht akzeptiert werden?
\item
Was macht die Turing-Maschine mit einem beliebigen Wort?
\end{teilaufgaben}

\begin{loesung}
\allowdisplaybreaks
\def\z#1{\hbox to0.4cm{\hfill\clap{$#1$}\hfill}}
\def\zb{\hbox to0.4cm{\hfill\clap{\blank}\hfill}}
\def\zt#1{\hbox to0.4cm{\hfill\clap{\texttt{#1}}\hfill}}
\begin{teilaufgaben}
\item
Es fehlen die Übergängen von $q_2$ und $q_3$ für das Bandzeichen \blank:
\begin{center}
\begin{tikzpicture}[>=latex,thick]
\coordinate (q2) at (-3,0);
\coordinate (q3) at (3,0);
\coordinate (qr) at (0,0);
\draw (q2) circle[radius=0.35];
\draw (q3) circle[radius=0.35];
\draw (qr) ellipse(0.65cm and 0.35cm);
\draw (qr) ellipse(0.60cm and 0.30cm);
\node at (q2) {$q_2$};
\node at (q3) {$q_3$};
\node at (qr) {$q_{\text{reject}}$};
\draw[->,shorten >= 0.65cm,shorten <= 0.35cm] (q2) -- (qr);
\draw[->,shorten >= 0.65cm,shorten <= 0.35cm] (q3) -- (qr);
\node at ($0.6*(q2)+0.4*(qr)$) [above] {$\scriptstyle\blank\blank,R$};
\node at ($0.6*(q3)+0.4*(qr)$) [above] {$\scriptstyle\blank\blank,R$};
\end{tikzpicture}
\end{center}
Sie sind im Zustandsdiagramm {\color{red}rot} ergänzt worden.
\item
Die Berechnungsgeschichte lautet:
\begin{align*}
&\zb   \zb     \z{q_0} \zt{1}  \zt{0}  \zt{1}  \zt{0}  \zb     \zb\\
&\zb   \zb     \zt{1}  \z{q_1} \zt{0}  \zt{1}  \zt{0}  \zb     \zb\\
&\zb   \zb     \z{q_2} \zt{1}  \zt{0}  \zt{1}  \zt{0}  \zb     \zb\\
&\zb   \zb     \zt{0}  \z{q_3} \zt{0}  \zt{1}  \zt{0}  \zb     \zb\\
&\zb   \zb     \z{q_4} \zt{0}  \zt{1}  \zt{1}  \zt{0}  \zb     \zb\\
&\zb   \z{q_4} \zb     \zt{0}  \zt{1}  \zt{1}  \zt{0}  \zb     \zb\\
&\zb   \zb     \z{q_0} \zt{0}  \zt{1}  \zt{1}  \zt{0}  \zb     \zb\\
&\zb   \zb     \zt{0}  \z{q_0} \zt{1}  \zt{1}  \zt{0}  \zb     \zb\\
&\zb   \zb     \zt{0}  \zt{1}  \z{q_1} \zt{1}  \zt{0}  \zb     \zb\\
&\zb   \zb     \zt{0}  \zt{1}  \zt{1}  \z{q_1} \zt{0}  \zb     \zb\\
&\zb   \zb     \zt{0}  \zt{1}  \z{q_2} \zt{1}  \zt{0}  \zb     \zb\\
&\zb   \zb     \zt{0}  \zt{1}  \zt{0}  \z{q_3} \zt{0}  \zb     \zb\\
&\zb   \zb     \zt{0}  \zt{1}  \z{q_4} \zt{0}  \zt{1}  \zb     \zb\\
&\zb   \zb     \zt{0}  \z{q_4} \zt{1}  \zt{0}  \zt{1}  \zb     \zb\\
&\zb   \zb     \z{q_4} \zt{0}  \zt{1}  \zt{0}  \zt{1}  \zb     \zb\\
&\zb   \z{q_4} \zb     \zt{0}  \zt{1}  \zt{0}  \zt{1}  \zb     \zb\\
&\zb   \zb     \z{q_0} \zt{0}  \zt{1}  \zt{0}  \zt{1}  \zb     \zb\\
&\zb   \zb     \zt{0}  \z{q_0} \zt{1}  \zt{0}  \zt{1}  \zb     \zb\\
&\zb   \zb     \zt{0}  \zt{1}  \z{q_1} \zt{0}  \zt{1}  \zb     \zb\\
&\zb   \zb     \zt{0}  \z{q_2} \zt{1}  \zt{0}  \zt{1}  \zb     \zb\\
&\zb   \zb     \zt{0}  \zt{0}  \z{q_3} \zt{0}  \zt{1}  \zb     \zb\\
&\zb   \zb     \zt{0}  \z{q_4} \zt{0}  \zt{1}  \zt{1}  \zb     \zb\\
&\zb   \zb     \z{q_4} \zt{0}  \zt{0}  \zt{1}  \zt{1}  \zb     \zb\\
&\zb   \z{q_4} \zb     \zt{0}  \zt{0}  \zt{1}  \zt{1}  \zb     \zb\\
&\zb   \zb     \z{q_0} \zt{0}  \zt{0}  \zt{1}  \zt{1}  \zb     \zb\\
&\zb   \zb     \zt{0}  \z{q_0} \zt{0}  \zt{1}  \zt{1}  \zb     \zb\\
&\zb   \zb     \zt{0}  \zt{0}  \z{q_0} \zt{1}  \zt{1}  \zb     \zb\\
&\zb   \zb     \zt{0}  \zt{0}  \zt{1}  \z{q_1} \zt{1}  \zb     \zb\\
&\zb   \zb     \zt{0}  \zt{0}  \zt{1}  \zt{1}  \z{q_1} \zb     \zb\\
&\zb   \zb     \zt{0}  \zt{0}  \zt{1}  \zt{1}  \zb     \z{q_a} \zb\\
\end{align*}
\item
Am Ende der Berechnung steht das Wort \texttt{0011} auf dem Band.
\item
Nein.
Die Übergänge nach $q_{\text{reject}}$ können gar nicht vorkommen.
Der Übergang von $q_2$ könnte nur stattfinden, wenn das Zeichen vor
der Null auch eine Null wäre, aber damit wäre die Maschine gar nicht
erst in den Zustand $q_1$ gekommen.
Der Übergang aus $q_3$ ist unmöglich, weil der Übergang nach $q_2$
sichergestellt hat, dass an der aktuellen Stelle eine \texttt{0}
sein muss.
\item
Die Turing-Maschine implementiert einen Sortieralgorithmus für die
Ziffern des Wortes.
\qedhere
\end{teilaufgaben}
\end{loesung}

\begin{bewertung}
Fehlende Übergänge ({\bf F}) 1 Punkt,
Berechnungsgeschichte ({\bf B}) 2 Punkte,
Resultatwort ({\bf R}) 1 Punkt,
Begründung dafür, dass alle Wörter akzeptiert werden ({\bf A}) 1 Punkt,
Sortieralgorithmus (gilt auch als Begründung für {\bf A}) ({\bf S}) 1 Punkt.
\end{bewertung}
